\section{Conclusion and Future Work}
In this paper, we proposed a robotic decision-making framework that integrates symbolic knowledge and uncertainty, with the goal of reducing user workload. The proposed method maintains a belief over possible worlds within a POMDP-based planning framework and quantitatively models uncertainty through a probability distribution over these worlds. Based on this, the robot can assess its own uncertainty and selectively query the human only when necessary. This approach enables stable task execution without continuous supervision and minimizes unnecessary interactions, thereby reducing user intervention.

In particular, we leverage entropy-based confidence to determine when to issue queries and select the most informative hypothesis to construct questions, improving the efficiency of interaction. Experimental results show that the proposed method reduces the number of human interactions while maintaining or improving task success rates, and effectively lowers user workload. User studies further support this trend, demonstrating that a system which actively manages uncertainty and requests assistance only when needed can operate effectively in real-world settings.

These results point to several extensions. First, richer feedback channels can build on the current binary hypothesis interface by allowing users to provide corrected symbolic facts or natural-language clarifications. Second, scalable frontier construction can extend the approach to larger domains through pruning, abstraction, or learned proposal distributions. Third, user-aware query selection can combine robot belief uncertainty with estimates of what users can judge from the interface. Finally, the framework can be extended to multi-robot settings, where multiple agents share uncertainty and coordinate when and whom to query. These extensions preserve the central idea of the framework: human feedback is most useful when it is requested at belief-relevant moments and expressed in a form that supports both planning and human understanding.
